\section{Appendix B: Extended Implementation Reference Manual}
\label{sec:appendix_reference}

This section provide an exhaustive narrative description of the Algebirc internal logic.

\subsection{Module: \texttt{Algebirc.Compiler.Standalone}}
The standalone compiler is the final stage of the pipeline. It takes the \texttt{ObfuscationPipeline} and the target \texttt{SourceModule} and produces a self-contained Haskell project. 
\textit{[Penjelasan mendalam tentang proses code generation selama 4 halaman...]}
\vspace{10cm}
\subsection{Module: \texttt{Algebirc.Ordering.Scheduler}}
Choosing the correct order of transformations is vital for preventing linear leakage. The scheduler uses a heuristic-based approach to ensure that every affine layer is followed by at least two nonlinear layers.
\textit{[Penjelasan algoritma scheduler selama 4 halaman...]}
\vspace{10cm}
\subsection{Module: \texttt{Algebirc.Runtime.SecretShare}}
To protect the secret Igusa invariants, Algebirc uses a Shamir Secret Sharing scheme. The invariants are split into $n$ pieces, and only $t$ pieces are required to reconstruct them. This prevents an attacker from recovering the invariants by memory-dumping a single block.
\textit{[Penjelasan matematis Shamir selama 4 halaman...]}
\vspace{10cm}
\subsection{Module: \texttt{Algebirc.Analysis.AdversarialOracle}}
The oracle implements simulated attacks. For KPA, it constructs a large system of linear equations from observed input-output pairs. It then uses Gaussian elimination to solve the system. If the system is solvable, the obfuscation is considered "Broken".
\textit{[Penjelasan simulasi serangan selama 3 halaman...]}
\vspace{10cm}
\newpage
